\documentclass[final,a4paper,12pt]{article}

\usepackage{notes}
\setlength\parskip{1ex}
\setlength\parindent{0cm}
\usepackage{mathtools}
\usepackage{amsthm}
\theoremstyle{definition}
\newtheorem{problem}{Problem}[section]
\newtheorem{solution}{Solution}[section]

\title{Geophysics III}
\author{Calculus Tutorial Solutions}


\begin{document}
\maketitle

\section{Derivative Practice Problems}

\begin{problem}
    Plate Cooling \\
    To find the heat flow from the temperature, we need to take the derivative
    of the temperature with respect to depth; the solution includes the
    following derivative
    \begin{equation*}
        \od{}{z} \frac{1}{n} \sin \left( \frac{n \pi z}{L} \right)
    \end{equation*}
\end{problem}

\begin{solution}
    Step 1:
    \begin{align*}
        \od{}{z} \frac{1}{n} \sin \left( \frac{n \pi z}{L} \right) 
        && \od{}{x} a u &= a \od{}{x} u \\
        && a &= \frac{1}{n} \\ 
        && u &= \sin \left( \frac{n \pi z}{L} \right) \\
        \frac{1}{n} \od{}{z} \sin \left( \frac{n \pi z}{L} \right)
    \end{align*}
    Step 2:
    \begin{align*}
        \frac{1}{n} \od{}{z} \sin \left( \frac{n \pi z}{L} \right)
        && \od{}{x} a u &= a \od{}{x} u \\
        && a &= \frac{1}{n} \\ 
        && u &= \sin \left( \frac{n \pi z}{L} \right) \\
        \frac{1}{n} \sin \left( \frac{n \pi z}{L} \right) \od{}{z} \left( \frac{n \pi z}{L} \right)
    \end{align*}
\end{solution}

\begin{problem}
    Seismic Reflection \\
    To find the minimum travel-time due to a seismic reflection off a boundary,
    we need to take the derivative of t(x) as below,
    \begin{equation*}
        \od{}{x} v_1^{-1} [ (h^2 + x^2)^{1/2} + (h^2 + (d - x)^2)^{1/2}]
    \end{equation*}
\end{problem}

\begin{problem}
    Decompression Melting \\
    To model partial melting, you need to take the derivative of the melt
    depression term,
    \begin{equation*}
        \od{}{F} K \left( \frac{X_0}{D + F(1 - D)} \right)^\gamma .
    \end{equation*}
\end{problem}

\begin{problem}
    Gravity on an Ellipsoid \\
    As part of the derivation to determine the gravity of an ellipsoid (remember
    we make an ellipsoidal correction to gravity), we need to differentiate the
    equation of an ellipsoid,
    \begin{equation}
    \od{}{\theta} \frac{1 - f}{1 - f\sin^2 \theta} .
    \end{equation}
\end{problem}




\section{Integral Practice Problems}

\begin{problem}
    Steady-State Temperature \\
    To find the temperature in a half-space, we need to integrate the thermal
    gradient with depth,
    \begin{equation*}
        \int \left( \frac{q_0}{k} - \frac{A}{k} z \right) \dd{z} ,
    \end{equation*}
    where $q_0$, $A$, and $k$ are constants.
\end{problem}

\begin{problem}
    Radioactive Decay \\
    To find the number of parent radionuclides present due to radioactive
    decay, we need to integrate both sides of the following equation,
    \begin{equation*}
        \int_{N_0}^N \frac{\dd{N'}}{N'} = - \int_0^t \lambda \dd{t'} ,
    \end{equation*}
    and once complete, solve for $N$.
\end{problem}

\begin{problem}
    Buried Slab \\
    This integral is similar to one we needed for our buried slab problem,
    \begin{equation*}
        \int_0^R \frac{r}{(r^2 + z^2)^{3/2}} \dd{r} .
    \end{equation*}
\end{problem}

\begin{problem}
    Mass Flux \\
    The mass flux of asthenosphere (mass passing a given point per unit time)
    moved by a tectonic plate can be expressed as
    \begin{equation*}
        \int_0^h \rho u_0 \frac{y}{h} \dd{y} .
    \end{equation*}
\end{problem}

\begin{problem}
    Entropy \\
    In petrology, the entropy is a key thermodynamic parameter.  The change in
    entropy is related to temperature as
    \begin{equation}
        \Delta S = \int_{T_1}^{T_2} \frac{C_P}{T} \dd{T} .
    \end{equation}
\end{problem}

\end{document}
